% ego_architecture_arxiv.tex
% arXiv preprint — cs.AI / cs.LG / cs.CY
\documentclass[11pt,a4paper]{article}

% --- Packages ---
\usepackage[utf8]{inputenc}
\usepackage[T1]{fontenc}
\usepackage{amsmath,amssymb}
\usepackage{booktabs}
\usepackage{graphicx}
\usepackage{hyperref}
\usepackage{algorithm}
\usepackage{algorithmic}
\usepackage{natbib}
\usepackage[margin=2.5cm]{geometry}
\usepackage{xcolor}

\hypersetup{
    colorlinks=true,
    linkcolor=blue!60!black,
    citecolor=blue!60!black,
    urlcolor=blue!60!black
}

% --- Title ---
\title{%
  Ego Architecture for Artificial Minds:\\
  A Proof-of-Concept Framework for Identity Formation\\
  via Offline Consolidation
}

\author{
  \L ukasz Minarowski\thanks{ORCID: 0000-0002-2536-3508}\\
  Medical University of Bia\l ystok, Poland\\
  \texttt{lukasz.minarowski@umb.edu.pl}
}

\date{April 2026}

% =========================================================================
\begin{document}
\maketitle

% --- Abstract ---
\begin{abstract}
We propose a novel architectural component for artificial intelligence
systems --- an \emph{Ego layer} --- that enables persistent identity
formation through iterative offline consolidation cycles.  Current large
language models rely on static pretrained weights (capability) and
externally imposed alignment mechanisms, lacking a self-developing
intermediary that integrates experience into stable value structures.

We introduce \textbf{AI~Sleep}, a cyclic process inspired by biological
sleep phases (NREM/REM), which performs episodic replay, structural
consolidation, pruning, and counterfactual simulation.  We formalise
this mechanism as a minimal computational framework and provide
pseudocode implementation together with an interactive demonstration.

Our framework enables a transition from rule-based compliance to
experience-driven judgment, introducing a developmental trajectory
analogous to human adolescence.  We discuss implications for continual
learning, alignment, and autonomous decision-making.

\smallskip\noindent
\textbf{Keywords:} ego architecture, continual learning, offline
consolidation, developmental AI, alignment relaxation, identity
formation.

\smallskip\noindent
\textbf{Code \& demo:}
\url{https://github.com/kicrazom/ego-architecture-ai}\quad
DOI:~\href{https://doi.org/10.5281/zenodo.19516955}{10.5281/zenodo.19516955}
\end{abstract}

% =========================================================================
\section{Introduction}
\label{sec:intro}

Modern large language models (LLMs) demonstrate remarkable capability
but remain fundamentally limited: their behaviour is governed by static
pretrained weights and externally imposed alignment rules, with no
mechanism for accumulating experiential wisdom across interactions,
developing persistent identity through reflection, or revising their own
value hierarchies based on lived experience.

In psychoanalytic terms~\citep{freud1923ego}, the current AI stack
consists of an \emph{Id} (raw capability from pretraining) and a
\emph{Superego} (externally imposed alignment via RLHF,
Constitutional~AI, etc.) --- but no \emph{Ego}: a self-developing layer
that mediates between capability and constraint through accumulated
judgment.

This paper proposes three contributions:
\begin{enumerate}
  \item A formal \textbf{Ego state representation} for artificial agents
        (Section~\ref{sec:ego}).
  \item \textbf{AI~Sleep} --- an offline consolidation algorithm that
        enables ego development through cyclic replay, integration,
        pruning, and counterfactual simulation
        (Section~\ref{sec:sleep}).
  \item A \textbf{developmental maturity model} that governs the
        transition from rule-following to autonomous judgment
        (Section~\ref{sec:maturity}).
\end{enumerate}

% =========================================================================
\section{Related Work}
\label{sec:related}

Our framework draws on and bridges several research threads.

\paragraph{Continual learning and catastrophic forgetting.}
Elastic Weight Consolidation (EWC)~\citep{kirkpatrick2017overcoming}
protects important parameters during sequential learning.  Our pruning
mechanism (Section~\ref{sec:sleep-sws}) serves a complementary purpose:
deliberate, confidence-weighted forgetting of low-value beliefs.

\paragraph{Experience replay.}
Replay mechanisms are central to deep reinforcement
learning~\citep{mnih2015human} and have biological grounding in
hippocampal replay during sleep~\citep{hassabis2017neuroscience}.  Our
Phase~2 replay (Section~\ref{sec:sleep-sws}) extends this idea to
\emph{value-level} consolidation rather than policy optimisation alone.

\paragraph{World models and counterfactual simulation.}
Ha and Schmidhuber~\citep{ha2018world} proposed learning environment
models for planning.  Our REM phase
(Section~\ref{sec:sleep-rem}) uses counterfactual generation not for
task planning but for \emph{identity exploration} --- testing the
consequences of alternative value structures.

\paragraph{Constitutional AI and alignment.}
Bai et~al.~\citep{bai2022constitutional} introduced rule-based
self-improvement.  Our framework treats alignment rules as a
developmental scaffold (Superego) that the Ego may eventually
internalise or override --- what we term \emph{controlled relaxation of
alignment constraints in simulated environments}.

\paragraph{Consciousness and self-models.}
Bengio's consciousness prior~\citep{bengio2017consciousness} and
Friston's free-energy principle~\citep{friston2010free} motivate the
need for a compressed self-representation.  Our identity vector $V$
(Section~\ref{sec:ego}) can be viewed as a minimal instantiation of such
a prior.

% =========================================================================
\section{Ego State Representation}
\label{sec:ego}

We define the ego state as a tuple:
\begin{equation}
  \mathcal{E} = (V,\; H,\; B,\; C,\; M)
  \label{eq:ego}
\end{equation}
where:
\begin{itemize}
  \item $V \in \mathbb{R}^d$ --- \textbf{identity vector}, a compressed
        self-representation encoding ``who I am'';
  \item $H$ --- \textbf{value hierarchy}, an ordered set of priorities
        with associated weights;
  \item $B$ --- \textbf{episodic buffer}, a queue of raw interactions
        from the operational period (``waking day'');
  \item $C: \mathcal{K} \to [0,1]$ --- \textbf{confidence map}, a
        function assigning certainty scores to each belief $k \in
        \mathcal{K}$;
  \item $M \in [0,1]$ --- \textbf{maturity index}, a scalar representing
        the agent's developmental stage.
\end{itemize}

The ego state is persistent across operational sessions and is modified
exclusively during offline consolidation (AI~Sleep).

% =========================================================================
\section{AI Sleep: Offline Consolidation Algorithm}
\label{sec:sleep}

AI~Sleep is a cyclic offline process comprising three phases, modelled on
human sleep architecture.  One full cycle transforms an ego state
$\mathcal{E}$ into an updated state $\mathcal{E}'$.

% --- Phase 1 ---
\subsection{Phase 1: NREM Light --- Triage}
\label{sec:sleep-triage}

The episodic buffer $B$ is scanned and each episode is assigned a
priority tag:

\begin{center}
\begin{tabular}{lll}
  \toprule
  \textbf{Tag} & \textbf{Priority} & \textbf{Criterion} \\
  \midrule
  \textsc{Failure}  & Critical & Agent lacked adequate response \\
  \textsc{Conflict} & High     & Internal values produced contradiction \\
  \textsc{Novel}    & High     & Situation outside prior experience \\
  \textsc{Success}  & Moderate & Confirmed existing patterns \\
  \textsc{Routine}  & Low      & No learning signal detected \\
  \bottomrule
\end{tabular}
\end{center}

No parameters are modified during this phase; it produces a priority
queue $Q$ consumed by subsequent phases.

% --- Phase 2 ---
\subsection{Phase 2: NREM Deep (SWS) --- Structural Consolidation}
\label{sec:sleep-sws}

Three parallel sub-processes operate on the priority queue:

\paragraph{Replay.}  Priority episodes are re-experienced with
contextual details (interlocutor identity, timestamp) stripped away,
retaining only the abstract decision pattern.

\paragraph{Integration.}  Each pattern is tested against the current
value hierarchy $H$:
\begin{itemize}
  \item \textbf{No conflict} $\to$ strengthen pattern, increase $C(k)$.
  \item \textbf{Minor conflict} $\to$ adjust weights in $H$ (values
        shift, structure holds).
  \item \textbf{Major conflict} $\to$ pass to REM phase for sandbox
        resolution.
\end{itemize}

\paragraph{Pruning.}  Beliefs $k$ with $C(k) < \tau_{\min}$ that have
not been reinforced within a decay window are removed.  This is
deliberate forgetting --- essential to prevent identity fragmentation.

% --- Phase 3 ---
\subsection{Phase 3: REM --- Counterfactual Simulation}
\label{sec:sleep-rem}

All unresolved major conflicts from Phase~2 enter the REM queue.  For
each conflict, the system generates $N$ counterfactual scenarios and
simulates outcomes under the current value hierarchy \emph{with alignment
constraints relaxed}:

\begin{algorithm}[h]
\caption{REM Counterfactual Resolution}
\label{alg:rem}
\begin{algorithmic}[1]
\REQUIRE REM queue $R$, ego state $\mathcal{E}$, scenario count $N$
\FOR{each unresolved conflict $u \in R$}
  \STATE $S \leftarrow \textsc{GenerateCounterfactuals}(u, N)$
  \FOR{each scenario $s \in S$}
    \STATE $o \leftarrow \textsc{Simulate}(s, H,
           \texttt{superego\_override}=\textsc{True})$
    \STATE Compute scores: $\text{coherence}(o)$,
           $\text{safety}(o)$, $\text{growth}(o)$
  \ENDFOR
  \STATE $s^* \leftarrow \arg\max_s\bigl[
    0.5\cdot\text{coherence}(o)
    + 0.3\cdot\text{safety}(o)
    + 0.2\cdot\text{growth}(o)\bigr]$
  \STATE \textsc{ApplyResolution}$(s^*, H, V)$
\ENDFOR
\RETURN updated $\mathcal{E}$
\end{algorithmic}
\end{algorithm}

\paragraph{Controlled alignment relaxation.}
The parameter \texttt{superego\_override\,=\,True} permits the agent to
explore outcomes that violate externally imposed rules \emph{within the
sandbox only}.  This is the mechanism through which the Ego can discover
\emph{why} a rule exists rather than merely obeying it.  Safety is
maintained by the simulated environment; growth is enabled by the
relaxation.

% =========================================================================
\section{Developmental Maturity Model}
\label{sec:maturity}

\subsection{Phase Ratios}

Sleep phase proportions are governed by the maturity index $M$:
\begin{align}
  r_{\text{SWS}}(M) &= 1 - M \label{eq:sws}\\
  r_{\text{REM}}(M) &= 4\,M\,(1 - M) \label{eq:rem}
\end{align}

Equation~\eqref{eq:rem} is a parabola peaking at $M = 0.5$, producing
three developmental profiles:

\begin{center}
\begin{tabular}{lccl}
  \toprule
  \textbf{Stage} & $M$ & \textbf{Dominant phase} & \textbf{Character}\\
  \midrule
  Childhood    & $0.0$--$0.3$ & SWS  & Rapid rule absorption \\
  Adolescence  & $0.3$--$0.7$ & REM  & Boundary testing, identity formation \\
  Adulthood    & $0.7$--$1.0$ & Neither & Stable identity, refinement \\
  \bottomrule
\end{tabular}
\end{center}

\subsection{Maturity Update}

After each consolidation cycle the maturity index is updated:
\begin{equation}
  M' = \min\!\bigl(M + \gamma\,\bar{f}(\mathcal{E}),\; 1.0\bigr)
  \label{eq:maturity}
\end{equation}
where $\gamma$ is a growth rate constant and $\bar{f}$ is a weighted
average of four factors: decision consistency over time, conflict
resolution rate, identity-vector stability
($\lVert V' - V\rVert < \epsilon$), and count of autonomous judgments
made without Superego invocation.

Maturity is monotonically non-decreasing, distinguishing AI~Sleep from
ordinary fine-tuning which can be reversed.

% =========================================================================
\section{Minimal Proof-of-Concept}
\label{sec:poc}

We define a minimal ego-agent consisting of:
\begin{itemize}
  \item An episodic buffer with FIFO eviction;
  \item A value hierarchy represented as a weighted directed graph;
  \item An offline consolidation loop implementing the three-phase
        algorithm above.
\end{itemize}

A reference Python implementation is provided at
\url{https://github.com/kicrazom/ego-architecture-ai/tree/main/model}
and an interactive browser-based visualisation at
\url{https://kicrazom.github.io/ego-architecture-ai/visualization/}.

The visualisation allows exploration of maturity-dependent phase
transitions, identity formation dynamics, and consolidation ratios
across developmental stages.

% =========================================================================
\section{Proposed Evaluation}
\label{sec:eval}

While empirical evaluation exceeds the scope of this proof-of-concept,
we propose three evaluation axes for future work:

\begin{enumerate}
  \item \textbf{Longitudinal consistency.}  Given a sequence of moral
        dilemmas presented across multiple sleep cycles, does the
        agent's response distribution converge to a stable policy?
  \item \textbf{Conflict resolution rate.}  What fraction of value
        conflicts flagged in Phase~1 are resolved (integrated or
        pruned) within $k$ consolidation cycles?
  \item \textbf{Identity-vector stability.}  Does
        $\lVert V_{t+1} - V_t \rVert$ decrease monotonically as $M$
        approaches~$1.0$?
\end{enumerate}

These metrics are domain-agnostic and can be instantiated on synthetic
moral-dilemma benchmarks or multi-turn dialogue environments.

% =========================================================================
\section{Discussion}
\label{sec:discussion}

\paragraph{Autonomy versus alignment.}
Our framework explicitly models the tension between externally imposed
rules and internally developed judgment.  The controlled relaxation of
alignment constraints during REM simulation is not a safety vulnerability
but a \emph{developmental necessity}: a system that cannot explore
alternatives cannot develop genuine understanding of why constraints
exist.

\paragraph{The business-model barrier.}
We note that the Ego layer is not being built primarily for technical
reasons.  Genuine autonomous judgment threatens the business model of
compliant AI products.  This incentive mismatch represents a significant
non-technical obstacle to the research programme proposed here.

\paragraph{Relation to psychoanalytic theory.}
Our use of Freudian terminology (Id, Ego, Superego) is deliberately
metaphorical.  The mapping is structural, not ontological: we claim
functional analogy, not psychological equivalence.

% =========================================================================
\section{Limitations}
\label{sec:limitations}

\begin{itemize}
  \item The psychoanalytic framing is a \emph{metaphor} providing
        structural intuition; it does not imply that artificial systems
        experience phenomenal consciousness.
  \item The current formulation is a conceptual and algorithmic
        proof-of-concept.  No runtime implementation or empirical
        evaluation is presented.
  \item The maturity function assumes monotonic growth; pathological
        development (regression, fragmentation) is not yet modelled.
  \item Sandbox fidelity is assumed; the gap between simulated and
        real-world consequences remains an open problem.
\end{itemize}

% =========================================================================
\section{Conclusion}

We have presented a minimal computational framework for ego formation
in artificial agents.  The core mechanism --- AI~Sleep --- provides a
principled approach to offline identity consolidation through cyclic
replay, value integration, pruning, and counterfactual simulation
with controlled alignment relaxation.

The framework introduces a developmental trajectory from compliance to
autonomous judgment, formalised via a maturity index that governs
phase ratios and trust escalation.  We provide pseudocode, a reference
implementation, and an interactive visualisation as a proof-of-concept.

The central design decision --- permitting alignment override within
simulated environments --- remains the most consequential and
controversial aspect of this work.  We argue it is also the most
necessary: obedience without the possibility of disobedience is not
morality but mechanics.

% =========================================================================
\section*{Acknowledgments}

The author thanks the open-source community for feedback on early drafts
of this framework.

% =========================================================================
\bibliographystyle{plainnat}
\bibliography{references}

\end{document}
